\section{Preprocessing}
\label{preprocessing}

The analysis process starts when \artiphishell is given a ``task'' from the competition organizers.
The task contains some metadata that uniquely identifies the task and states the acceptable time frame for analysis, as well as several archives, such as the OSS-Fuzz bundle for the target application, a baseline repository, and, in the case of a delta-mode task, the diff containing the new code that has been introduced)

\subsection{\builder}
The first step after receiving the task is performed by the \builder agent.
\builder creates various artifacts, depending on the language of the target application.
\begin{description}
\item[Canonical build:] the compiled application in its standard format.
\item[Debug build:] a version of the application with debug symbols.
\item[\afluzzer build:] a version of the application with the instrumentation required by AFL++.
\item[\jazzmine build:] a version of the application with the instrumentation required by Jazzer.
\item[Coverage build:] a version of the application with coverage specific information that is used to produce detailed coverage reports when a specific seed is sent to a harness.
\end{description}

\subsection{\analyzer}
After \builder creates the artifacts previously described, the \analyzer agent extracts metadata about the target and stores it in the \analysisgraph database.

The following are some of the type of data that are extracted:
\begin{description}
\item[WLLVMbear:] contains all the compiler commands executed during the building of the various artifacts.
\item[CodeQL:] contains the CodeQL database for the target, and the results of running a set of queries to extract the control-flow graph (CFG), the data-flow graph (DFG), and the code-property graph (CPG). 
\item[Seed Corpus:] an initial pre-determined corpus of input seeds based on popular formats and protocols.
\end{description}
The \analyzer agent is also responsible for the identification of the endpoints that represent harnesses and will be used in the vulnerability analysis process.

\subsection{\indexer}
After the \analyzer has completed its tasks, the \indexer agent creates a separate repository of function-related metadata.
This repository is critical for the efficient and consistent access to function-related information across all agents.
This is achieved using CLANG's and ANTLR's indexing capabilities to produce a list of functions with a unique identifier and a clear association between functions and their locations in the code base.
This information is exposed using a dedicated API with caching to improve the speed and scalability of access.







